\documentclass[11pt]{article}
\usepackage[T1]{fontenc}
\usepackage[utf8]{inputenc}
\usepackage{lmodern}
\usepackage[margin=1in]{geometry}
\usepackage{amsmath,amssymb}
\usepackage{graphicx}
\usepackage{booktabs}
\usepackage{array}
\usepackage{tabularx}
\usepackage{multirow}
\usepackage{xcolor}
\usepackage[numbers,sort&compress]{natbib}
\usepackage{microtype}
\usepackage[colorlinks=true,linkcolor=blue!50!black,citecolor=blue!50!black,urlcolor=blue!50!black]{hyperref}

\newcommand{\gap}{\mathrm{gap}}
\newcommand{\tbd}[1]{\textcolor{red}{[#1]}}
\newcommand{\BT}{Base$\leftrightarrow$Turbo}
\newcolumntype{L}{>{\raggedright\arraybackslash}X}
\graphicspath{{figures/}}

\title{What Recurs and What Is Inherited:\\
Layerwise Spectral Telemetry of Generative Transformers\\[4pt]
\large Mode-I: an empirical atlas and its falsification record}
\author{Zac Doyle\\ \small mode:LIMINAL\\ \small \url{https://github.com/Z-L-D/mode-I}}
\date{October 2026}

\begin{document}
\maketitle

\begin{abstract}
We record three spectral statistics of the token activation matrix (a normalised amplitude participation ratio $g$, the leading singular-value gap and the spectral entropy $H$) at every layer and every sampling step of five diffusion and flow transformers (Z-Image Turbo and Base, Mochi~1, Stable Diffusion~3.5 Medium and Large), with language-model controls, and ask which features of the resulting depth--time trajectories recur across architectures and which are inherited under distillation. The answer depends on the observable. Scalar trajectories of $g$ look alike across unrelated models; on linear $g$ the best unrelated pair beats the distillation pair by $0.049$ in correlation. As the comparison becomes stricter (log scale, level field, registered depth shape, depth tangent), the margin of the distillation pair over the best unrelated pair grows to $+0.047$, $+0.131$, $+0.152$ and $+0.703$: depth tangents agree at $0.94$--$0.99$ between Z-Image Base and its distilled Turbo and at most $0.24$ for any unrelated pair. Prompt-relational geometry is shared broadly (Mantel $\rho=0.13$--$0.43$ between unrelated models, $0.56$--$0.62$ within the lineage), and the ordered depth profile at the first sampling state decodes ten semantic domains at $0.43$--$0.61$ (chance $0.10$). Along the sampling clock, Z-Image MLP shows expansion, contraction and re-expansion of layer-mean $g$; along depth, the layer step of $g$ follows a $+,-,+,-$ sign sequence in every image and video architecture at some noise level, beating marginal and mean-field surrogates in some channels and not in others. Two out-of-family checks, a convolutional U-Net and an audio flow transformer (ACE-Step~1.5, ten preregistered predictions), keep the low-dimensional backbone and a depth step that is not a function of $g$, but not the strong anisotropy; the audio model does not show the depth sign sequence. We report with equal weight the results that did not survive: a mid-run synchronisation peak, a universal semantic clock, sharp seed commitment, an autoregressive scaling anomaly and an apparent architecture ordering of scalar curves. Every claim is scoped to its capture regime, and the preregistration and correction trail is released with the code.
\end{abstract}

\tableofcontents
\newpage

%=====================================================================
\section{Introduction}\label{sec:intro}

Comparisons between generative models usually look at their outputs or at static representations of fixed inputs \citep{raghu2017svcca,kornblith2019similarity}. A diffusion or flow transformer, however, computes a trajectory: at each sampling step the network transforms its input through a stack of layers, and the sampler then moves the latent to a lower noise level \citep{ho2020denoising,song2021scorebased,lipman2023flow,liu2023flow}. The object we study is the evolution of layerwise spectral geometry along both axes of that computation, depth and sampling time.

We use a fixed passive instrument (CoS-PAFS) that records, for every layer, step and generation, three statistics of the token activation matrix: the normalised amplitude participation ratio $g$, the spectral gap and the spectral entropy. Five diffusion and flow transformers were captured on a common set of prompts. Two of them form a distillation lineage: Z-Image Turbo is a few-step model distilled within the Z-Image family (from its SFT model, with distribution-matching distillation and reward post-training), and Z-Image Base is the undistilled model of the same family \citep{zimageteam2025zimage,liu2025decoupled,tongyi2026zimage}. The others are unrelated architectures or a scaling pair. Two further models were captured afterwards as out-of-family checks: the convolutional U-Net of Stable Diffusion~1.5 (\S\ref{sec:unet}) and an audio flow transformer, ACE-Step~1.5 (\S\ref{sec:ace}).

We ask empirical questions and do not assume a universal answer:
\begin{enumerate}
 \item Is the measured trajectory low-dimensional, and is that more than smoothness?
 \item Do trajectories of unrelated models resemble each other, and at which level of description?
 \item Does depth geometry carry information about prompt content, and is the content geometry shared across models?
 \item Is there a recurrent sign structure along the sampling clock or along depth, and does it survive null models?
 \item What can passive telemetry say about when a seed's outcome becomes predictable?
 \item Which interventions (guidance, ablation, scale, distillation) preserve which geometry?
\end{enumerate}

\paragraph{Release boundary.} This paper makes no claim about a separatrix or commitment boundary in the sampling dynamics, about any relation between the geometry and output quality, about per-token phases in autoregressive models, or about a renormalisation-group universality class. The data cannot support these claims, and Section~\ref{sec:limits} lists the experiments that could.

\paragraph{Contributions.}
\begin{enumerate}
 \item A multi-model telemetry atlas with named capture regimes and a released reproducer for every number (\S\ref{sec:instrument}).
 \item Low-dimensional depth--time organisation of the measured geometry, beyond smoothness (\S\ref{sec:backbone}).
 \item An observable hierarchy: scalar resemblance crosses architecture boundaries, while registered shape and especially depth tangents separate the distillation lineage (\S\ref{sec:similarity}).
 \item Semantic information in ordered depth geometry and broadly shared prompt-relational geometry (\S\ref{sec:content}).
 \item Two separately scoped sign-structure results, along the sampling clock and along depth, with bootstrap and surrogate qualification (\S\ref{sec:topology}).
 \item Calibrated estimators for seed predictability that remove a sample-size artefact (\S\ref{sec:seed}), an intervention hierarchy (\S\ref{sec:interventions}), and a falsification record that defines the claim boundary (\S\ref{sec:falsification}).
 \item Preregistered out-of-family checks on a convolutional U-Net (\S\ref{sec:unet}) and on an audio flow transformer (\S\ref{sec:ace}), which separate the results that hold outside the image-transformer setting from those that do not.
\end{enumerate}
The mathematical treatment of these results (a pooling theorem, projection degeneracy, the non-autonomous flow family and null derivations) is given in the companion paper, Mode-II.

%=====================================================================
\section{Related work}\label{sec:related}

\paragraph{Representation geometry along depth.} Linear and nonlinear dimension estimates of hidden representations change systematically with depth \citep{ansuini2019intrinsic,recanatesi2019dimensionality}, and their coarse depth profiles recur across very different transformers \citep{valeriani2023geometry,cheng2025emergence}. A recurring coarse shape is therefore expected and is not by itself evidence of shared computation. Our observable is closest to the participation ratio of neuroscience \citep{gao2017theory,litwinkumar2017optimal} and to the entropy effective rank \citep{roy2007effective,garrido2023rankme}; finite-sample bias of such ratios \citep{chun2026estimating} is why every number here is reported with its sample size. In a diffusion model, intrinsic dimension has been measured across prompts and layers \citep{kvinge2023exploring} and depth$\times$timestep roles of DiT blocks have been described \citep{chen2024deltadit}; we add the joint depth--time trajectory of the spectrum across several model families.

\paragraph{Comparing models.} Representational similarity is usually compared with CCA-type measures, CKA or RSA \citep{raghu2017svcca,morcos2018insights,kornblith2019similarity,kriegeskorte2008representational}. These measures disagree, have known pathologies and need calibration \citep{ding2021grounding,davari2023reliability,cui2022deconfounded,klabunde2025similarity}, and apparent cross-model convergence \citep{huh2024platonic} can shrink under permutation calibration \citep{groger2026revisiting}. Shape metrics \citep{williams2021generalized} formalise the registered-shape level we use. Our contribution here is empirical: on the same data, envelope, field, shape, content and tangent comparisons rank model pairs differently.

\paragraph{Distillation.} Few-step students are obtained by trajectory matching \citep{salimans2022progressive,song2023consistency}, adversarial objectives \citep{sauer2024adversarial,sauer2024fast} or distribution matching \citep{yin2024onestep,yin2024improved}, the family used for Z-Image Turbo \citep{liu2025decoupled}. Students can keep teacher concepts and pairings \citep{gandikota2026distilling,surkov2025onestep,li2026why} or restructure internal computation \citep{haskins2025distilled}, and recent language-model distillation aligns depth trajectories rather than single layers \citep{chi2026mta}. We measure, without access to the training objective, how much of the depth dynamics survives.

\paragraph{Regimes of the sampling dynamics.} Statistical-physics analyses describe symmetry breaking, speciation and collapse along the reverse process \citep{raya2023spontaneous,biroli2023generative,biroli2024dynamical,sclocchi2025phase}, and critical windows in which features are decided \citep{li2024critical}; recent work reads class commitment from conditional entropy \citep{handke2026entropic}. Sampling trajectories of ODE samplers share a low-dimensional regular shape \citep{chen2024trajectory}. These works measure outputs or clone continuations. Our telemetry is passive and internal, which is why it can describe predictability and persistence but not commitment (\S\ref{sec:seed}).

\paragraph{Nulls for trending series.} Trending series correlate spuriously \citep{yule1926why,granger1974spurious}, and circular-shift or phase-randomised surrogates can be beaten by drift alone \citep{theiler1992testing,schreiber2000surrogate,harris2020nonsense}. This literature motivates the trend-preserving, mean-field and marginal surrogates used throughout.

%=====================================================================
\section{Instrument, data and provenance}\label{sec:instrument}

\subsection{Observables}
At each sampling step and each transformer block, a forward hook takes the token activation matrix $X\in\mathbb{R}^{N\times D}$ of one submodule (the attention output or the MLP output) and computes the singular values $s_1\ge\dots\ge s_D$ of $X/\sqrt N$, i.e.\ the square roots of the eigenvalues of the uncentred second-moment matrix $X^\top X/N$. It records
\begin{equation}
 g=\frac{1}{D}\frac{(\sum_is_i)^2}{\sum_is_i^2},\qquad \gap=\frac{s_1}{s_2},\qquad H=-\sum_iq_i\log q_i,\ q_i=\frac{s_i}{\sum_js_j}.
\end{equation}
$Dg$ is the amplitude participation ratio, the squared ratio of nuclear to Frobenius norm; it is related to, but not the same as, the eigenvalue participation ratio \citep{gao2017theory} and the effective rank $\exp(H)$ of \citet{roy2007effective}. No mean is subtracted. Normalised depth is $u=(\ell-1)/(L-1)$.

\subsection{Capture regimes}
With classifier-free guidance \citep{ho2022classifierfree} every step evaluates a conditional and an unconditional branch. The hook can record the statistic of the stacked branches (\emph{pooled}), of the conditional branch only (\emph{cond-only}), or the mean of the two branch statistics (\emph{branch-mean}). These are different functionals of the same computation. Pooling provably raises $g$ above the energy-weighted branch mean (Mode-II, pooling theorem), and on identical SD3.5 Medium trajectories it raises $\log g$ by a median of $+0.040$, $+0.050$ and $+0.088$ at $1024$, $768$ and $512$ pixels. The excess follows how far the branches diverge, not matrix size: per (generation, step, layer) cell it predicts one minus the logged prompt influence with held-out $R^2=0.62$--$0.64$ at each resolution, and once that divergence is modelled, resolution adds only $\Delta R^2=0.011$, whereas an isotropic baseline with $N$ versus $2N$ rows would move $\log g$ by $0.05$--$0.38$ over the same sizes (preregistered, Round 8 R8-4). The regime differs across the archive (Table~\ref{tab:corpus}) and, for SD3.5, within campaigns: SD3.5 Medium switches regime inside runs and SD3.5 Large varies by domain. Every number in this paper names its regime, and absolute $g$ is never compared across regimes.

\subsection{Models and campaigns}
\begin{table}[!t]
\centering
\small
\caption{Telemetry campaigns in the archive. $L$: blocks (for the U-Net, its transformer blocks); $D$: hidden width; Gens: complete generations. Regime: capture functional (\S\ref{sec:instrument}). All Z-Image and SD3.5 captures use sampling shift $3.0$; the Mochi shift and frame size are not recorded.}
\label{tab:corpus}
\begin{tabularx}{\linewidth}{lcccccL}
\toprule
Model & $L$ & $D$ & Steps & CFG & Gens & Hooks; capture regime\\
\midrule
Z-Image Turbo (ZIT) & 30 & 3840 & 8 & 1.0 & 480 & attn, MLP; cond-only\\
Z-Image Base (ZIB) & 30 & 3840 & 25 & 4.0 & 476 & attn, MLP; pooled\\
Mochi 1 & 48 & 3072 & 30 & 4.5 & 478 & attn (visual stream); branch-mean\\
Mochi 1 & 48 & 3072 & 100 & 1.0 & 60 & attn; cond-only\\
SD3.5 Medium & 24 & 1536 & 40 & 4.5 & 480 & attn (image stream); pooled at step 0, branch-mean from step 2\\
SD3.5 Large & 38 & 2432 & 28 & 3.5 & 485 & attn (image stream); pooled or branch-mean by domain\\
\midrule
ZIB, Rounds 3--4 & 30 & 3840 & 25 & 1, 4 & 19--60 & attn, MLP; cond-only or explicit branches\\
SD3.5 M, Rounds 3--4 & 24 & 1536 & 40 & 1, 4.5 & 30--60 & attn; $512/768/1024$ px; explicit branches\\
SD1.5 U-Net (R3) & 16 & 320--1280 & 50 & 7.5 & 476 & attn, MLP; $512$ px; explicit branches\\
ACE-Step 1.5 turbo (R3) & 24 & 2048 & 6 & 1.0 & 120 & attn, MLP; audio, 180 s; cond-only\\
\midrule
Qwen2.5 1.5B/3B/7B & 28/36/28 & 1536--3584 & 256 tok & -- & 120 each & one hook; running statistic\\
Mistral 7B v0.3 Instr./ablated & 32 & 4096 & 256 tok & -- & 120 each & one hook; running statistic\\
\bottomrule
\end{tabularx}
\end{table}

The diffusion campaigns use a common pool of 120 prompts in ten researcher-defined semantic domains, about four seeds per prompt; 117 prompts are shared by Z-Image Base, Turbo and Mochi. Z-Image hooks read the attention output and the second MLP projection, including the 64 text rows of the joint sequence; SD3.5 and Mochi hooks read the image or visual stream only. Images are $1024$ pixels except where stated. Z-Image \citep{zimageteam2025zimage}, SD3.5 \citep{esser2024scaling} and Mochi~1 \citep{genmo2024mochi} are rectified-flow or flow-matching transformers \citep{liu2023flow,lipman2023flow,peebles2023scalable}. Turbo was run with its recommended few-step, guidance-free settings \citep{tongyi2025zimageturbo}. All Z-Image and SD3.5 captures used the ComfyUI \citep{comfyui2023} sampling shift of $3.0$, verified by fitting the logged $\sigma$ schedule; for Z-Image Base this is below the official value of $6.0$, and it is the setting under which Base and Turbo share a $\sigma$ axis. The language-model captures \citep{qwen2024qwen25,jiang2023mistral} record a running statistic over a growing context and are used only as fixed-weight and scaling controls (\S\ref{sec:interventions}).

\begin{figure}[t]
 \centering
 \includegraphics[width=\linewidth]{I-F1_corpus_instrument_map}
 \caption{Corpus and capture instrument. Each row is one campaign. The strip shows which capture functional the stored value records along normalized sampling progress. ComfyUI's CFG batching decides this, so it differs between campaigns and can change within a run: cond-only; pooled cond\,+\,uncond; branch mean; pooled or branch mean by prompt domain; explicit per-branch rows (instrument v2); for the autoregressive models, a nested running statistic (the v1 LLM script re-encodes the full context at every token). For the SD1.5 U-Net (\S\ref{sec:unet}), layers are its 16 transformer blocks. The ZIB archive is pooled throughout; ZIT is cond-only; SD3.5\,M is pooled at step 0, mixed at step 1 and branch-mean from step 2; SD3.5\,L is pooled in 74\% of steps, depending on domain. Columns give steps, transformer layers, hooked submodule, CFG scale and generations/prompts. Every number in this paper names its capture regime.}
 \label{fig:corpus}
\end{figure}

\subsection{Coordinates}
For diffusion models the sampling coordinate is the noise level $\sigma$. Z-Image Turbo and Base are compared at equal $\sigma$: the optimal affine time warp between them improves registered correlation by at most $0.005$ over equal-$\sigma$ matching, and per prompt equal-$\sigma$ matching loses at most $0.01$ in 99\% (attention) and 100\% (MLP) of prompts. Every one of 117 shared prompts selects zero layer shift. Models of different depth are compared on normalised depth $u$. Because equal $\sigma$ is not equal signal-to-noise across resolutions \citep{hoogeboom2023simple}, resolution comparisons use the SD3 resolution shift \citep{esser2024scaling}.

The autoregressive captures are not per-token trajectories: in all seven language-model channels the recorded statistic is perfectly rank-correlated with context length ($|\rho|=1.000$), because it is accumulated over the growing context. We do not use them for any per-token dynamical claim.

\subsection{Reproducibility and preregistration}
The archive was analysed in seven audit rounds. From Round 2 onward, each round froze its predictions with a SHA-256 hash and a UTC time before any calculation \citep{nosek2018preregistration,simmons2011false}; preregistered and post hoc results are labelled throughout. Rounds 5 and 6 scored $13/9$ and $17/5$ (confirmed/refuted). Round 7 was a release audit of existing outputs; its eight checks are consistency checks, not new predictions. Appendix~\ref{app:prereg} lists all scorecards, and every refuted prediction is reported.

%=====================================================================
\section{A low-dimensional backbone}\label{sec:backbone}

For each model we stack the population depth profiles of $(\log g,\log\gap,H)$, standardised across depth at each sampling state, into a states$\times$features matrix and compute its principal components. PC1 explains $0.71$--$0.90$ of the variance and the effective rank is $1.5$--$2.5$ (Fig.~\ref{fig:backbone}); the Round 1 reproduction with the entropy effective rank gives $0.74$--$0.89$ and $1.53$--$2.37$, and definition variants move these by about $0.05$. Turbo has eight states, so its rank is at most seven by construction. As a null, every time slice's depth profile is independently circularly shifted along depth, which keeps each profile's smoothness but breaks its registration across steps; all five models beat all 999 null campaigns ($p=0.001$), whose mean effective rank is $4.9$--$14.8$. The low rank therefore comes from depth profiles that stay registered across sampling steps, not from smoothness of individual profiles. The SD1.5 U-Net (\S\ref{sec:unet}, post hoc) has the same backbone: PC1 $0.90$ (attention) and $0.62$ (MLP), effective rank $1.46$ and $2.24$, against null ranks of $11.1$ and $12.1$ ($p=0.001$). The audio model ACE-Step (\S\ref{sec:ace}, post hoc) gives PC1 $0.78$ and $0.56$ and effective rank $2.10$ and $2.73$ against null ranks of $3.9$ and $4.1$ ($p=0.001$ and $0.011$); with six sampled states its rank is at most five.

The same holds within every semantic domain: across 68 model$\times$submodule$\times$domain trajectories PC1 is at least $0.70$, and model and submodule identity explain $0.90$--$0.96$ of the variation in rank and slowdown summaries against $0.02$--$0.05$ for domain. This is low dimensionality of the measured spectral representation; it says nothing about the dimension of the hidden-state dynamics.

Along the sampling clock, every channel shows an early high-speed transient at $\sigma\approx1$ followed by a $\sigma$-invariant plateau. The plateau onset, defined as the first $\sigma$ at which population state speed falls below twice its median below $\sigma=0.6$, lies at $\sigma=0.78$--$0.87$ in all eight archive channels (four architectures, CFG 1--4.5, 8--100 steps), and is tighter in $\sigma$ than in step fraction (spread $0.32$--$0.64$). On the clean SD3.5 Medium conditional-branch captures it moves with resolution, from $0.72$ at $512$ to $0.81$ at $1024$; on the resolution-corrected clock the three resolutions give $0.81$--$0.84$. The SD1.5 U-Net, mapped to the same clock at equal signal-to-noise ratio, reaches its plateau later, at $0.58$ (attention) and $0.67$ (MLP), or $0.73$ and $0.80$ after the resolution shift (post hoc; the comparison depends on that mapping, \S\ref{sec:unet}). The audio model ACE-Step, on its own flow-matching $\sigma$, reaches it at $0.90$ (attention) and $0.80$ (MLP), inside the preregistered window $[0.73,0.92]$ (B8, \S\ref{sec:ace}); with six steps, its plateau speed rests on a single interval. At matched $\sigma$, geometry changes by at most $2$--$3\%$ between 8 or 40 and 100 sampling steps (Turbo: $g$ within $1.9\%$, on a four-generation 100-step control; SD3.5 Medium: $3\%$ at CFG 1 and no change at CFG 4.5).

\begin{figure}[t]
 \centering
 \includegraphics[width=\linewidth]{I-F2_backbone}
 \caption{Low-dimensional temporal backbone in CoS spectral coordinates (descriptive). \textbf{Left:} fraction of variance on the first principal component of the population trajectory, after within-state depth standardization; \textbf{Right:} effective rank; \textbf{Grey markers:} mean of 999 circular-shift nulls, in which each time slice is shifted in depth independently. PC1 is 0.71--0.90 against null means of 0.19--0.40, and effective rank is 1.5--2.5 against 4.9--14.8 ($p = 0.001$ for every model). ZIT has only 8 sampled states, so its rank is at most 7 by construction. \textbf{Below the dotted line} (post hoc): the SD1.5 U-Net, conditional branch, attention (circle) and MLP (square), with the same definition and null: PC1 0.90 and 0.62, rank 1.46 and 2.24 against null ranks 11.1 and 12.1 ($p=0.001$). Diamonds: the audio model ACE-Step (attention filled, MLP open; 6 states): PC1 0.78 and 0.56, rank 2.10 and 2.73 against 3.9 and 4.1 ($p=0.001$, $0.011$).}
 \label{fig:backbone}
\end{figure}

%=====================================================================
\section{Similarity depends on the observable}\label{sec:similarity}

\subsection{Scalar envelopes}
Early in the project, layer-mean attention $g$ along $\sigma$ correlated at $r=0.9405$ between Z-Image Base and Mochi, slightly above Base and its own distillation Turbo ($0.9325$), with Turbo$\leftrightarrow$Mochi at $0.8217$. These values reproduce exactly from the archive. They are not evidence about architecture. Each campaign has a different capture regime (Base pooled, Turbo cond-only). On $\log g$ the same archive gives Base$\leftrightarrow$Turbo $0.973$ against Base$\leftrightarrow$Mochi $0.860$, and with Base's conditional branch captured explicitly (Round 4, 19 prompts) the linear-$g$ values are $0.989$ and $0.831$ (Fig.~\ref{fig:scalar}). More fundamentally, every denoising trajectory shares a coarse trend in $\sigma$, and a scalar correlation is bounded below by that shared trend whatever the residual dynamics \citep{yule1926why} (Mode-II, trend-dominance proposition). On linear $g$ the best unrelated pair (Turbo$\leftrightarrow$SD3.5 Large, $0.981$) exceeds the lineage pair ($0.933$).

\begin{figure}[t]
 \centering
 \includegraphics[width=\linewidth]{I-F4_capture_regime_reversal}
 \caption{A scalar envelope ordering that reverses with the capture regime (worked example). \textbf{(a)} Layer-mean attention $g(\sigma)$, each model scaled to its own maximum. Every envelope rises, so envelope correlations are dominated by a shared monotone trend. \textbf{(b, c)} Correlation of the ZIB envelope with ZIT (filled) and with Mochi (open) when only the ZIB capture is substituted, on linear $g$ (b) and log $g$ (c). The historical ordering ZIB$\leftrightarrow$Mochi 0.941 $>$ ZIB$\leftrightarrow$ZIT 0.933 appears only for the pooled archive on linear $g$. It is already reversed on the same 19 prompts. It is reversed in every explicit capture (cond, uncond, branch mean, CFG\,1), and on log $g$ in every capture (archive 0.973 vs 0.860; CFG\,1 cond-only 0.9998 vs 0.732). This is an example of capture-regime sensitivity, not evidence about architecture.}
 \label{fig:scalar}
\end{figure}

\subsection{Fields, shapes and tangents}
We compare model pairs at matched $\sigma$ at four further levels: $\log g$ envelopes, the level field over (depth, $\sigma$), the depth shape (profiles standardised across depth at each state, correlated over a common 24-point depth grid), and the depth tangent (cosine of the first differences of those profiles along depth). Table~\ref{tab:hierarchy} gives the lineage margin, the \BT{} value minus the best unrelated pair, at each level (preregistered, Round 6 S1).

\begin{table}[h]
\centering\small
\caption{Lineage margin by observable level.}
\label{tab:hierarchy}
\begin{tabular}{lccccc}
\toprule
 & linear $g$ & $\log g$ & field & shape & tangent\\
\midrule
Base$\leftrightarrow$Turbo & 0.933 & 0.973 & 0.984 & 0.981 & 0.942\\
Best unrelated pair & 0.981 & 0.926 & 0.853 & 0.829 & 0.240\\
Margin & $-0.049$ & $+0.047$ & $+0.131$ & $+0.152$ & $+0.703$\\
\bottomrule
\end{tabular}
\end{table}

Registered shape at matched $\sigma$ is $0.981$ (attention) and $0.983$ (MLP) for \BT{} on the pooled archive, and $0.9997$ on the clean Base CFG 1 capture; unrelated Z-Image$\leftrightarrow$Mochi shape still reaches $0.79$--$0.83$, so shape alone does not isolate the lineage. The tangent does. For \BT{} it depends on how Base was captured, in a fixed order: $0.94$/$0.97$ (attention/MLP) against the pooled archive, $0.981$/$0.989$ against the unconditional branch, $0.988$/$0.991$ against the branch mean, $0.992$/$0.992$ against the conditional branch at CFG 4, and $0.999$ against Base run at CFG 1. Guidance therefore moves conditional geometry by less than $1\%$, and most of the archive's shortfall from $1$ is the pooled capture (guidance moves the tangent by $0.007$, the capture by $0.025$--$0.049$). In attention that shortfall sits in the deepest third of the normalised depth grid ($0.86$, against $0.97$ and $0.96$ in the first two thirds), and every third is at least $0.98$ on the conditional capture. The \BT{} tangent does not depend on sample size: on 50 prompt-matched subsets of 60 generations it lies at $0.937$--$0.946$ (attention) and $0.965$--$0.968$ (MLP) against $0.942$ and $0.967$ on all generations (preregistered, Round 8 R8-5, R8-6, R8-13). Every unrelated pair has $|\cos|\le0.24$, inside a layer-permutation null (95th percentile $0.24$/$0.20$). The scaling pair SD3.5 Medium$\leftrightarrow$Large gives $-0.25$. The tangent is unchanged by any rotation of the hooked activations and by embedding them in a wider space, so a width difference cannot by itself lower it; by the same token, a high tangent is necessary for an orthogonal correspondence between two models' activations but does not establish one (Mode-II, \S4.3).

\begin{figure}[t]
 \centering
 \includegraphics[width=\linewidth]{I-F3_observable_hierarchy}
 \caption{Lineage becomes identifiable only at the level of the local depth derivative. \textbf{(a)} Similarity of ZIT$\leftrightarrow$ZIB (Base$\leftrightarrow$Turbo lineage), SD3.5\,M$\leftrightarrow$L (same family) and the eight unrelated pairs, at five observable levels (attention): L0: scalar envelope, linear and log $g$; L1: population field; L2: registered depth shape; L3: matched-$\sigma$ tangent cosine. On the scalar envelope an unrelated pair, ZIT$\leftrightarrow$SD3.5\,L (0.981), beats the lineage (0.933). The tangent separates the lineage: 0.94 against $\le 0.24$ for every unrelated pair. \textbf{(b)} Base$\leftrightarrow$Turbo tangent cosine under every ZIB capture functional, attention and MLP. Values run from 0.943/0.968 (archive, pooled) through 0.981/0.989 (CFG\,4 uncond), 0.988/0.991 (branch mean) and 0.992/0.992 (CFG\,4 cond) to 0.999/0.999 (CFG\,1 cond-only). All lie far above the best unrelated pair (ZIB$\leftrightarrow$Mochi, 0.24).}
 \label{fig:hierarchy}
\end{figure}

\subsection{Aligning states does not align dynamics}
Can unrelated models be made to agree by aligning their state spaces? Procrustes, linear and CCA maps fitted on state landmarks, without derivative labels, visibly improve state alignment between unrelated population fields. Their matched depth derivatives remain near zero or anti-aligned: the best map for each unrelated pair averages a derivative correlation of $0.007$, while \BT{} averages $0.937$ across maps and channels (Fig.~\ref{fig:d10}).

\begin{figure}[t]
 \centering
 \includegraphics[width=\linewidth]{I-F9_state_vs_dynamics}
 \caption{Aligned states do not imply aligned local dynamics. \textbf{(a)} Registered depth-shape correlation (state) against matched-$\sigma$ tangent cosine (dynamics) for all 10 model pairs (attention). ZIT$\leftrightarrow$Mochi and ZIB$\leftrightarrow$Mochi have state similarity 0.79--0.83 but tangent cosines of only 0.21--0.24. Only Base$\leftrightarrow$Turbo is high on both (0.98, 0.94). \textbf{(b)} State-only alignment maps (raw, Procrustes, linear, CCA), learned from state landmarks without derivative labels. They raise the matched state correlation of unrelated and same-family fields (up to 0.75), while the matched depth-derivative correlation stays near zero ($-0.18$ to $0.09$). Base$\leftrightarrow$Turbo keeps a derivative correlation of 0.92--0.95 under every map.}
 \label{fig:d10}
\end{figure}

%=====================================================================
\section{Content geometry}\label{sec:content}

\subsection{Semantic domains are decodable from ordered depth geometry}
A multinomial decoder given only the ordered layerwise profile of $(g,\gap,H)$ at the first measured state recovers the ten prompt domains at $0.45$--$0.61$ with prompt-grouped cross-validation (chance $0.10$; recomputed for this paper with the Round 1 decoder, which gave $0.43$--$0.61$). Layer-averaged scalars give $0.11$--$0.29$ and a prompt-length control $0.15$--$0.18$; a TF-IDF decoder on the prompt text itself reaches $0.38$--$0.42$ (Table~\ref{tab:decode}). Scrambling the layer order within each generation, which keeps every per-layer value, destroys the signal: the real decoder beats all 99 scramble realisations in every channel ($p=0.01$; scramble maximum $\le0.19$). Because the SD3.5 capture regime varies by domain, the decoder could in principle read the regime. The SD1.5 U-Net (post hoc) also beats every scramble, but only at the level of the text decoder: $0.38$ (attention) and $0.42$ (MLP) against TF-IDF $0.40$ (Table~\ref{tab:decode}, Fig.~\ref{fig:decode}). For the audio model ACE-Step the test is not informative: its 29 music prompts give only two or three prompts per domain, so with prompt-grouped folds even the text decoder is at chance ($0.10$), and the telemetry gives $0.11$ (attention) and $0.12$ (MLP). Within regime-homogeneous subsets decoding remains $4.0$--$4.7\times$ chance for pooled SD3.5 Large (9 domains, $0.45$--$0.55$) and $2.3$--$3.4\times$ chance for branch-mean subsets ($0.58$--$0.82$ on 4 SD3.5 Large domains; $0.30$--$0.34$ on 10 SD3.5 Medium domains). Across the seven transformer channels, first-state decoding tracks the first-state prompt variance share (Spearman $\rho=0.64$, stable to leaving out any channel; preregistered, Round 9 R9-A4).

\begin{table}[h]
\centering\small
\caption{Ten-way domain decoding at the first state (archive regimes; chance $0.10$). Scramble: maximum over 99 layer-order scrambles. Below the line: the SD1.5 U-Net (\S\ref{sec:unet}; post hoc, conditional branch, 476 generations, 119 prompts), same decoder and features. ACE-Step (\S\ref{sec:ace}; post hoc, cond-only, 120 generations, 29 prompts, 2--3 per domain) is not informative: every decoder, including TF-IDF, is near chance.}
\label{tab:decode}
\begin{tabular}{lccccc}
\toprule
Channel & Ordered profile & Layer-averaged & Prompt length & TF-IDF on text & Scramble max\\
\midrule
ZIT attn & 0.548 & 0.285 & 0.175 & 0.410 & 0.167\\
ZIT MLP & 0.554 & 0.190 & 0.175 & 0.410 & 0.135\\
ZIB attn & 0.489 & 0.191 & 0.168 & 0.406 & 0.139\\
ZIB MLP & 0.460 & 0.113 & 0.168 & 0.406 & 0.135\\
Mochi attn & 0.609 & 0.283 & 0.159 & 0.412 & 0.140\\
SD3.5 M attn & 0.448 & 0.185 & 0.150 & 0.417 & 0.167\\
SD3.5 L attn & 0.548 & 0.243 & 0.157 & 0.382 & 0.188\\
\midrule
SD1.5 attn & 0.379 & 0.187 & 0.134 & 0.395 & 0.172\\
SD1.5 MLP & 0.424 & 0.183 & 0.134 & 0.395 & 0.187\\
ACE-Step attn & 0.108 & 0.033 & 0.100 & 0.100 & 0.108\\
ACE-Step MLP & 0.117 & 0.067 & 0.100 & 0.100 & 0.175\\
\bottomrule
\end{tabular}
\end{table}

The domains are prompt categories in text-conditioned models. Decoding them shows that the conditioning signal is written into ordered depth geometry; it is not spontaneous class formation of the kind studied in unconditional diffusion \citep{biroli2024dynamical,sclocchi2025phase}.

\begin{figure}[t]
 \centering
 \includegraphics[width=\linewidth]{I-F5_semantic_decoding}
 \caption{Ten-way semantic-domain decoding from telemetry. Decoder: logistic regression on $\log|[g, \mathrm{gap}, H]|$ per layer, 5-fold cross-validation grouped by prompt; chance 0.10. \textbf{(a)} At the first measured state the ordered depth profile reaches 0.45--0.61 across seven channels. This beats: the layer-mean scalars (0.11--0.29); prompt length (0.15--0.18); a TF-IDF decoder on the prompt text itself (0.38--0.42); every one of 99 depth-scrambled realizations (layer order permuted per generation; maximum 0.19, $p = 0.01$). \textbf{(b)} Along sampling the ordered profile stays at 0.34--0.61, above the layer mean in every channel and at every step. Black (post hoc): the SD1.5 U-Net, conditional branch, on 119 prompts. Its ordered profile gives 0.38 (attention) and 0.42 (MLP) at the first state, about the TF-IDF level, and 0.30--0.42 along sampling, still above its layer mean (0.14--0.19). Diamonds in (a), right of the second line: ACE-Step (audio; 29 prompts, 2--3 per domain), where every decoder including TF-IDF is near chance, so the test is uninformative.}
 \label{fig:decode}
\end{figure}

\subsection{Prompt-relational geometry is shared across models}
For each model and sampling state we form the matrix of distances between prompt-mean depth profiles and compare models by representational similarity analysis \citep{kriegeskorte2008representational} with Mantel permutation tests \citep{mantel1967detection}. The prompt-relational structure is shared across all five diffusion models (preregistered, Round 6 S2): Mantel $\rho=0.13$--$0.43$ for unrelated pairs and $0.56$--$0.62$ for \BT{}, every $p\le0.0005$. It persists within themes (29 of 30 cells), after partialling out prompt-text overlap, and after partialling out prompt length (partial Mantel on character-length differences: partial over full $\rho\ge0.998$ in all 24 unrelated-pair cells, every $p\le0.0001$; preregistered, Round 8 R8-12). Shared text encoders do not explain it: Z-Image$\leftrightarrow$Mochi ($0.39$) exceeds SD3.5 Medium$\leftrightarrow$Large ($0.31$), which share encoders (a refuted prediction, S2e). The SD1.5 U-Net shares it too (post hoc, bottom row of Fig.~\ref{fig:rsa}): on the 116 prompts common to all six models its attention gives $\rho=0.11$--$0.35$ against the five transformers, every $p=0.0005$. The language models share it too, at matched context (preregistered, Round 9 R9-B6 and its text-control addendum). The prompt-mean depth profiles are compared at the prompt and at generated tokens 128 and 255, residualised on log context length, with the Mantel test partialled on context-length and TF-IDF text distance. Mistral 7B Instruct and Llama~3.1 8B Instruct, from different developers with different tokenizers and training data, give $\rho=0.44$--$0.48$ on 118 shared prompts, and the three Qwen2.5 sizes give $0.72$--$0.89$ on 24 prompts, every $p=0.0001$. The text controls remove at most $3\%$ of the effect, as they do for the diffusion pairs (unrelated pairs at the first state $0.23$--$0.38$ under the same two controls, at most $5\%$ removed). Each model's own agreement with the text distance is only $0.11$--$0.23$. Because a language model's measured tokens include the prompt itself, the reading is that unrelated models organise the same text alike, beyond its wording and length.

\begin{figure}[t]
 \centering
 \includegraphics[width=\linewidth]{I-F6_prompt_rsa}
 \caption{Prompt-similarity structure is shared across unrelated models and is strongest in the lineage. \textbf{(a--c)} Mantel $\rho$ between the $117\times117$ prompt-dissimilarity matrices of every model pair at the first, middle and last sampling step. The lineage cell is outlined. Base$\leftrightarrow$Turbo: 0.56--0.62; Unrelated pairs: 0.13--0.43; Same family (SD3.5\,M$\leftrightarrow$L): 0.31--0.43; Every cell has $p = 0.0005$ against its own 2000-permutation null (null p95 $\approx 0.02$). Bottom row, below the line (post hoc): SD1.5 U-Net attention against each transformer on the 116 prompts shared by all six models (0.11--0.35, every $p=0.0005$); the transformer cells use 117 prompts. \textbf{(d)} Robustness (post hoc). Full RSA against within-theme RSA (theme means removed, theme-preserving null) and against partial Mantel controlling for prompt-text RSA, for all 10 pairs $\times$ 3 steps. The shared structure persists within themes and after the prompt-text overlap is removed.}
 \label{fig:rsa}
\end{figure}

\subsection{Variance shares}
Across the archive, the system$\times$time scaffold accounts for $96\%$ of normalised depth-geometry variance and theme, prompt and seed for the remaining $4\%$. That share is mostly between models and between sampling states. Within a fixed (model, state) cell, prompt identity explains $51\%$--$96\%$ of the variance as a naive share and $36\%$--$95\%$ after correcting for seed replication; the lowest values are Base at $\sigma=1$, where the pooled first step carries a large seed share. The SD1.5 U-Net falls inside both ranges ($0.67$--$0.75$ naive, $0.56$--$0.67$ corrected; post hoc). The audio model ACE-Step lies at the top ($0.86$--$0.97$ naive, $0.82$--$0.96$ corrected); its share is highest at $\sigma=1$ and dips at $\sigma=0.60$, so it does not rise first as in every image channel (B10, refuted; \S\ref{sec:ace}). Both statements are true, and the second is the relevant one for content.

\subsection{Theme residuals}
After removing each model's scaffold, the alignment of theme residuals between models is strongest in the lineage at every state ($0.678$, $0.750$, $0.626$ at the first, middle and last states), but unrelated pairs reach $0.48$ by the middle state and $0.52$ at the end, while at the first state they lie at $0.07$--$0.22$. Theme inheritance is lineage-specific only at the first state (Fig.~\ref{fig:S1}). Across the ten pairs with both measurements, content similarity and tangent similarity are not detectably rank-coupled (Spearman $0.297$, $p=0.41$; a low-powered test).

%=====================================================================
\section{Sign structure along two axes}\label{sec:topology}

The archive supports two different three-segment statements. They are measured on different axes and are reported separately.

\subsection{Along the sampling clock: Z-Image MLP}
Layer-mean MLP $g$ in Z-Image follows expansion, contraction (by about $50\%$) and re-expansion as $\sigma$ decreases. A two-breakpoint piecewise-linear fit \citep{muggeo2003estimating} recovers the order $+,-,+$ in 1000 of 1000 generation bootstraps, on the archive and on clean explicit-branch captures at CFG 1 and CFG 4 (preregistered, S4). For the 25-step Base channel no time-order permutation fits comparably (0 of 1000 in all three captures); Turbo's eight states are too few for that test (38 of 1000). Attention $g$ instead rises and then falls late, $+,+,-$, in all five models (post hoc). The clock-axis statement is therefore specific to Z-Image MLP, the only transformer with MLP hooks. The SD1.5 U-Net (\S\ref{sec:unet}) has MLP hooks, and neither of its envelopes follows either pattern: with the same fit on the matched clock $t=\sigma/(1+\sigma)$, MLP and attention $g$ both rise over the whole schedule, $+,+,+$ in 1000 of 1000 prompt bootstraps (the MLP levels off and ends $0.6\%$ below its maximum), on raw $\sigma$ and on the resolution-shifted clock as well (post hoc; its schedule covers $t=0.94$--$0.12$). The time-averaged MLP depth profile peaks at layer 5, has a floor in the middle layers and rises over the last layers (post hoc, descriptive). The audio model ACE-Step (\S\ref{sec:ace}, post hoc) shows neither pattern either: on its six steps both envelopes peak at $\sigma=1$, drop by $16\%$ (MLP) and $19\%$ (attention) at the first step, partly recover and then fall, $-,+,-$ in 1000 of 1000 prompt bootstraps. Six points barely constrain a three-segment fit, so this is descriptive.

\begin{figure}[t]
 \centering
 \includegraphics[width=\linewidth]{I-F7_sigma_three_regime}
 \caption{Along $\sigma$, Z-Image MLP expands, contracts and re-expands; attention does not. \textbf{(a)} Population-mean layer-mean MLP $g(\sigma)$, scaled to its own peak, with the continuous three-segment piecewise-linear fit, for four captures: archive ZIT (8 steps); archive ZIB (pooled, CFG\,4); the clean ZIB CFG\,4 cond branch; ZIB CFG\,1 cond-only. Every channel gives $+\,-\,+$ in 100\% of 1000 prompt bootstraps, with a contraction of about 50\% (50--55\%). In the 25-step ZIB channels 0/1000 time-order permutations fit comparably; ZIT's 8 points cannot carry that test (38/1000). \textbf{(b)} The same fit on attention in all five models gives $+\,+\,-$ (a rise followed by a late fall), in 92--100\% of bootstraps, not $+\,-\,+$. Black, dashed (post hoc): the SD1.5 U-Net, conditional branch, MLP in (a) and attention in (b), drawn on the matched clock $t=\sigma/(1+\sigma)$. Both rise over the whole schedule ($+\,+\,+$ in 100\% of bootstraps), with no contraction and no late fall. Diamonds, dash-dot (post hoc): the audio model ACE-Step (6 steps, cond-only), MLP in (a) and attention in (b): $-\,+\,-$ in 100\% of bootstraps, with the maximum at $\sigma=1$.}
 \label{fig:sigma}
\end{figure}

\subsection{Along depth: the layer step of $g$}
The CoS-PAFS analyzer measures the layer-direction flow as $\beta=g(\ell+1)-g(\ell)$, bins all (generation, step, layer) points of $\beta$ against $g$ into 29 bins with at least 10 points each, and reports the sign pattern and zero crossings of the binned mean. We ran the analyzer's own function, unchanged. The three-crossing pattern $+,-,+,-$ recurs across architectures (Table~\ref{tab:beta}): in all 100 generation bootstraps for SD3.5 Medium and Large attention, in $90\%$ and $84\%$ for Z-Image Turbo and Base MLP, and in $67\%$ for Base attention. At a single noise level it appears in every image and video architecture at some stage of sampling: SD3.5 Large in 24 of 28 steps ($\sigma$ from $0.96$ to $0.19$), Mochi in 18 of 30, SD3.5 Medium in 13 of 40 (early, $\sigma\ge0.86$), Base MLP in 9 of 25 and Turbo MLP in 3 of 8.

Two surrogates qualify it. Surrogate A permutes each layer's $g$ across all rows (keeps per-layer marginals); Surrogate B permutes $g$ across generations within each (layer, step) cell (keeps the mean field). In SD3.5 Medium attention and Z-Image Base attention and MLP neither surrogate reproduces the pattern (0 of 30 each), so it requires within-generation coupling across layers. In Turbo MLP it beats Surrogate A but follows from the mean field (B reproduces it 30 of 30). In SD3.5 Large both surrogates reproduce it: it is null-generic. Mochi at CFG 4.5 and Turbo attention give a single crossing in the aggregate. So does the audio model ACE-Step, where the pattern was preregistered and is refuted (B6, \S\ref{sec:ace}): attention has a single crossing at every one of its six steps and in 99 of 100 prompt bootstraps, and surrogates reproduce that single crossing; its MLP shows three crossings at 2 of 6 steps.

\begin{table}[h]
\centering\small
\caption{Layer-axis $\beta$ sign pattern (analyzer, archive). Boot.: share of 100 generation bootstraps with $+,-,+,-$. Steps: share of single-$\sigma$ clouds with $+,-,+,-$. A, B: surrogate draws (of 30) reproducing the observed aggregate pattern. Below the line: the audio model ACE-Step (\S\ref{sec:ace}; 6 steps, cond-only; the attention row is preregistered B6, the MLP row post hoc).}
\label{tab:beta}
\begin{tabular}{llccccc}
\toprule
Channel & Aggregate & Boot. & Steps & A & B & Qualification\\
\midrule
SD3.5 M attn & $+-+-$ & 100\% & 13/40 & 0 & 0 & beats both\\
SD3.5 L attn & $+-+-$ & 100\% & 24/28 & 30 & 30 & null-generic\\
ZIT MLP & $+-+-$ & 90\% & 3/8 & 0 & 30 & beats A only\\
ZIB MLP & $+-+-$ & 84\% & 9/25 & 0 & 0 & beats both\\
ZIB attn & $+-+-$ & 67\% & 8/25 & 0 & 0 & beats both\\
Mochi attn (CFG 4.5) & $+-$ & 19\% & 18/30 & 30 & 7 & one crossing\\
ZIT attn & $+-$ & 5\% & 2/8 & 30 & 30 & one crossing\\
\midrule
ACE-Step attn & $+-$ & 0\% & 0/6 & 30 & 30 & one crossing\\
ACE-Step MLP & $+-$ & 5\% & 2/6 & 30 & 30 & one crossing\\
\bottomrule
\end{tabular}
\end{table}

No universal crossing location or interval ratio exists: the share of the middle segment ranges from $3\%$ (SD3.5 Medium) to $77\%$ (Base MLP), and the crossing count depends on the capture regime, on whether attention and MLP are pooled, on which steps enter the cloud, and on sample size. Within one channel the crossing also moves with $\sigma$: per step, Base MLP's first $+\to-$ crossing lies at higher $g^*$ inside its clock-axis contraction window ($\sigma=0.82$--$0.97$, above) than outside it (Mann--Whitney $p=0.004$; $p=0.011$ after scaling by each step's median $g$, post hoc), and the per-step three-crossing pattern occurs mostly on the plateau (preregistered, Round 8 R8-7). The robust statement is the sign sequence with its bootstrap rate and surrogate class, per channel and per $\sigma$. Figure~\ref{fig:beta}c shows the binned curves themselves, read back from the analyzer's plot (data and extraction script in the repository).

\begin{figure}[tp]
 \centering
 \includegraphics[width=\linewidth]{I-F8_layer_beta_topology}
 \caption{Along the layer axis, the analyzer's layer-direction flow $\beta_{\text{layer}} = g(L{+}1) - g(L)$ shows the $+\,-\,+\,-$ three-crossing pattern in Z-Image, SD3.5\,M/L and Mochi at some or all $\sigma$ (post hoc analyzer replay). \textbf{(a)} Crossing pattern at each denoising step, one $\sigma$ at a time. \textbf{(b)} Fraction of 100 generation bootstraps with $+\,-\,+\,-$ on all steps pooled: SD3.5\,M and L attention: 100\%; ZIT MLP: 90\%; ZIB MLP: 84\%; ZIB attention: 67\%; Mochi and ZIT attention: 19\% and 5\%. The panel also shows how often marginal (Surrogate A) and mean-field (Surrogate B) surrogates reproduce the observed pattern (30 draws). The pattern exceeds both surrogates in SD3.5\,M attention and in ZIB attention and MLP; SD3.5\,L is null-generic. \textbf{(c)} The binned curve $\beta_{\text{layer}}(g)$ itself for each channel, all steps pooled, as drawn by the analyzer's own \texttt{plot\_beta\_layer} (29 bins, at least 10 points per bin; blue: binned mean; grey band: $\pm1$ s.d.; dashed lines: the analyzer's crossings $g^{*}$, orange $+\to-$, green $-\to+$). In several channels the middle segment is narrow compared with the spread across generations. Read the pattern per $\sigma$ or with its bootstrap rate, not as a single crossing count.}
 \label{fig:beta}
\end{figure}

%=====================================================================
\section{Seed predictability without commitment language}\label{sec:seed}

For each sampling step we predict a generation's final seed residual (its state with the prompt mean removed) from its current residual, with ridge regression and prompt-grouped cross-validation, and record the $R^2$ curve $A(t)$.

The usual onset, the $\sigma$ at which $A(t)$ first exceeds $0.5$ ($\sigma_{50}$), turned out to be mostly a sample-size number. The ceiling $A(T-1)$ rises from $0.48$ to $0.89$ as the number of generations grows from 60 to 480, and $\sigma_{50}$ moves with it, from undefined (in 98\% of partitions at $n=60$) to $0.72$ and then $0.94$. An earlier contrast between ``early'' onsets in mixed-regime arms and ``late'' onsets in conditional arms was a comparison across different $n$. Cross-validated estimates at this sample size carry large error bars \citep{varoquaux2018cross}. We therefore use the relative onset, the $\sigma$ at which $A(t)/A(T-1)$ crosses $1/2$, reported at matched $n$ over prompt-grouped partitions. Across channels its range is $0.277$, against $0.501$ for $\sigma_{50}$, and its interval is narrower in 7 of 7 channels (preregistered, S3; Table~\ref{tab:seed}). Relative onsets lie at $\sigma=0.84$--$0.95$ (Turbo MLP $0.67$), and guidance shifts them by at most $+0.03$.

\begin{table}[h]
\centering\small
\caption{Seed predictability at matched $n$ (60 prompts, 50 prompt-grouped subsets, archive regimes): median and 5--95\% interval. VR4/walk: four-step variance ratio divided by its simulated finite-walk value.}
\label{tab:seed}
\begin{tabular}{lccc}
\toprule
Channel & Relative onset & $\sigma_{50}$ & VR4/walk\\
\midrule
ZIT attn & 0.841 [0.798, 0.867] & 0.514 [0.396, 0.608] & 1.13\\
ZIT MLP & 0.673 [0.627, 0.707] & 0.413 [0.362, 0.459] & 0.65\\
ZIB attn & 0.934 [0.904, 0.948] & 0.871 [0.802, 0.918] & 0.54\\
ZIB MLP & 0.887 [0.804, 0.939] & 0.754 [0.611, 0.854] & 0.49\\
Mochi attn & 0.902 [0.888, 0.910] & 0.862 [0.840, 0.879] & 0.85\\
SD3.5 M attn & 0.950 [0.928, 0.962] & 0.913 [0.865, 0.934] & 0.61\\
SD3.5 L attn & 0.930 [0.916, 0.946] & 0.867 [0.812, 0.910] & 0.87\\
\bottomrule
\end{tabular}
\end{table}

Seed deviations evolve gradually. Their four-step variance ratio \citep{lo1988stock}, calibrated, because the small-sample statistic is biased below one \citep{kim1991mean,lo1989size}, against simulated finite random walks of the same length (whose raw $\mathrm{VR}(4)\approx0.82$), lies at $0.49$--$0.87$ of the walk value, between an AR(0.8) process and a random walk; Turbo attention is the exception at $1.13$, and that value is what its 8-state grid produces: subsampling the 25--40-step channels to the 8 states nearest Turbo's schedule lifts VR4/walk above one in 3 of 5 (to $1.13$--$1.35$), while their relative onsets move by at most $0.023$ (preregistered, Round 8 R8-2). Guidance contracts conditional-branch seed spread: on SD3.5 Medium at $1024$ pixels with shared seeds, the across-seed spread at CFG 4.5 is $0.637$ of that at CFG 1 (prompt bootstrap $0.45$--$0.87$). We make no statement about the timing of that contraction.

These are statements about gradual predictability and persistence of seed deviations under passive observation. They are not a committor or fate boundary, which would require cloned continuations from fixed intermediate states \citep{e2006towards,bolhuis2002transition,e2010transitionpath}.

\begin{figure}[t]
 \centering
 \includegraphics[width=\linewidth]{I-F10_estimators}
 \caption{Seed predictability rises gradually, and the raw $\sigma_{50}$ estimator depends on sample size. \textbf{(a)} Archive SD3.5\,M: raw $\sigma_{50}$ against the ceiling-free relative onset as the number of generations $n$ grows. $\sigma_{50}$ moves from undefined in 98\% of resamples at $n = 60$ to 0.94 at $n = 480$; the relative onset stays at 0.87--0.96. \textbf{(b)} At matched $n = 60$ (50 subsets; median and 5--95\% range), the relative onset is 0.84--0.95 in $\sigma$ in every channel except ZIT MLP (0.67). \textbf{(c)} VR(4) on the calibrated scale: observed divided by a simulated random walk, with AR(0.8) ($\approx 0.85$--$0.88$) and iid-jitter ($\approx 0.30$) references. Seed deviations lie between AR(0.8) and a random walk (0.49--0.87 of walk), except ZIT attention (1.13), a value its 8-state grid produces (subsampling other channels to 8 states gives 1.13--1.35; Round 8). $\sigma_{50}$ values are not compared across different $n$, and no commitment step is implied.}
 \label{fig:seed}
\end{figure}

%=====================================================================
\section{An intervention hierarchy}\label{sec:interventions}

Different interventions move the geometry by very different amounts (Fig.~\ref{fig:interventions}).
\begin{itemize}
 \item \textbf{Fixed weights, changed inference.} Mochi with CFG 1.0 versus 4.5 keeps registered shape at $0.9965$ (also a 100- versus 30-step comparison). Abliteration of Mistral 7B \citep{arditi2024refusal} keeps shape at $0.99977$; at prefill the median $|\Delta\log g|$ is $0.0013$, layers 0--5 are bit-identical, and the final layer moves four times more than any other. The depth tangent is kept: at prefill the tangent cosine between the two models is $0.9998$ (prompt bootstrap $0.99976$--$0.99983$; Round 9 R9-B3). After the prompt the sampled texts diverge at a median of generated token 7, so only the prefill is a clean comparison. SD3.5 Medium at $512$, $768$ and $1024$ pixels on the clean Round 4 captures (CFG 4.5, explicit branches, 30 prompts) moves the main $\beta$ crossing only within $g^*=0.074$--$0.085$, and the pooled Round 3 captures at the same three resolutions all show the depth sign sequence $+,-,+,-$ (24-bin estimator of the review rounds); the conditional plateau onset moves from $0.72$ to $0.76$ to $0.81$, and onto $0.84$, $0.81$, $0.81$ on the resolution-corrected clock. Resolution therefore shifts the operating point in $\sigma$ rather than the organisation. (A legacy single-prompt suite gives shape $0.989$--$0.996$ across resolutions; Appendix~\ref{app:legacy}.)
 \item \textbf{Scaling.} Qwen2.5 1.5B, 3B and 7B keep broad state resemblance (registered correlation $0.45$--$0.63$) but not depth derivatives: at matched $n$ the derivative correlation is $+0.63$ for 1.5B$\leftrightarrow$3B (prompt bootstrap $0.61$--$0.65$, above a layer-permutation null whose 95th percentile is $0.38$), $-0.12$ for 1.5B$\leftrightarrow$7B ($-0.16$ to $-0.08$) and $-0.06$ for 3B$\leftrightarrow$7B ($-0.10$ to $-0.01$), and in all three pairs it rises monotonically with context length (preregistered, Round 9 R9-B2). Scaling can therefore keep part of the local derivative, but not reliably. The three sizes share one chat-template regime (identical prefill lengths for all 24 prompts). SD3.5 Medium$\leftrightarrow$Large fails registration throughout, with shape $0.32$ and tangent $-0.25$.
 \item \textbf{Distillation.} Z-Image Base$\to$Turbo preserves registered shape ($0.97$--$0.99$) and tangents ($0.94$--$0.99$), and the geometry-derived registration closes the MLP local eigenmode sequence ($p=0.002$ over all $8!$ time orders) but not attention ($p=0.26$).
 \item \textbf{Capture functional.} With the trajectory held fixed, pooling the branches shifts $\log g$ by $+0.04$ to $+0.09$ (\S\ref{sec:instrument}).
\end{itemize}
Guidance, ablation and capture move the operating point or deform shape slightly. Scaling preserves coarse resemblance but not local dynamics. Distillation preserves geometry and, especially, depth tangents.

\begin{figure}[t]
 \centering
 \includegraphics[width=\linewidth]{I-F11_intervention_hierarchy}
 \caption{Depth-shape deformation orders interventions. \textbf{(a)} Normalized shape deformation $1-r$: fixed-weight comparisons: $2\times10^{-4}$ to $1.1\times10^{-2}$ (Mistral Instruct$\leftrightarrow$abliterated 0.9998; Mochi CFG\,1$\leftrightarrow$4.5 0.9965; SD3.5\,M resolution pairs); distillation (ZIB$\leftrightarrow$ZIT): 0.027; capacity scaling: 0.37--0.77. The SD3.5 resolution rows (open markers) use one prompt at CFG\,1 with the legacy instrument. The Mochi CFG comparison is also 30 vs 100 steps. \textbf{(b)} Local-dynamics alignment where it was measured: matched-$\sigma$ tangent cosine for diffusion (Mochi CFG 0.99; ZIB$\leftrightarrow$ZIT 0.94/0.97 attn/MLP; SD3.5\,M$\leftrightarrow$L $-0.25$), the prefill depth-tangent cosine for the Mistral abliteration (0.9998, Round 9), and median depth-derivative $r$ at matched context length for Qwen (0.63, $-0.12$, $-0.06$). The Qwen prompt-bootstrap 5--95\% intervals are narrower than the markers; gray ticks mark each pair's layer-permutation null 95th percentile (0.38--0.45), which only 1.5B$\leftrightarrow$3B exceeds. \textbf{(c)} For fixed-weight comparisons, shape deformation against the absolute operating-point shift $A$ (log-RMS, 0.001--0.22). Fixed-weight changes move the operating point while leaving the shape almost unchanged.}
 \label{fig:interventions}
\end{figure}

%=====================================================================
\section{Falsification results that define the claim boundary}\label{sec:falsification}

The following results were proposed during the project and did not survive. They are part of the empirical result.

\paragraph{A mid-run synchronisation peak.} A cross-layer synchronisation index $R(\tau)$ peaks at $0.81$--$0.93$ near mid-sampling in every channel, and the values reproduce exactly. The peak is not an event. A straight-line surrogate with no dynamics gives a peak at least as high in 7 of 7 channels (preregistered, S5), only 3 of 7 channels beat a monotone-warp null at $p<0.05$, and centring-free co-movement measures have no interior maximum in 7 of 7. What survives is optional and supplementary: cross-layer timing coordination beyond independent monotone drift in Turbo attention and MLP ($p=0.014$, $0.008$) and, marginally, Base attention ($p=0.048$).

\paragraph{A universal semantic clock.} The step of maximal semantic Fisher information ranges from $0.14$ to $0.69$ in normalised time across channels; there is no common semantic peak.

\paragraph{Universal curvature, susceptibility and seed collapse.} Passive susceptibility peaks sit at the endpoints or are mutually incompatible; seed-share collapse is model-dependent and flat after the initial transient, and an elevated seed share at $\sigma\approx1$ coincides with pooled capture.

\paragraph{An autoregressive scaling anomaly.} A multiscale $\beta$ ladder in Qwen2.5-7B survived two nulls, but the recorded language-model statistic is a running statistic over context (rank correlation $1.000$ with context length in all seven channels), so the ladder is a context-clock effect and is demoted (Fig.~\ref{fig:S3}). The perfect rank correlation is a property of the instrument, not of the models: a nested sample of i.i.d.\ rows with no model at all, re-measured as rows are added, gives $|\rho(n,g)|=|\rho(n,gD/n)|=|\rho(n,H)|=1$ for isotropic, power-law and spiked populations at $D=1536$ and $4096$ (preregistered, Round 9 R9-B1). Every v1 capture has $n\le389<D$, so $g\le n/D$, but $g$ sits at $0.33$--$0.37$ of that ceiling; the ceiling alone does not produce the effect.

\paragraph{Architecture ordering of scalar curves.} The historical ZIB$\leftrightarrow$Mochi $>$ ZIB$\leftrightarrow$ZIT ordering is a capture-regime and projection effect (\S\ref{sec:similarity}).

\paragraph{Pooled-cloud fixed points.} Crossings of the pooled $(g,\beta)$ cloud were once read as fixed points. Under marginal and mean-field surrogates some are null-generic (\S\ref{sec:topology}), and dense multi-crossings are the temporal profile folded into $(g,\beta)$.

\paragraph{A rotation index.} The raw cosine between successive layer displacements clusters at $-1/2$, which is exactly what independent layers give (Mode-II, differencing-null proposition); after correction it adds no held-out information beyond depth and $\sigma$.

\begin{figure}[t]
 \centering
 \includegraphics[width=\linewidth]{I-F12_sync_falsification}
 \caption{The historical synchronization peak is not a mid-run event. \textbf{(a)} Cross-layer order parameter $R(\sigma)$ of the population-mean field: observed against a per-layer straight-line surrogate (ZIB and SD3.5\,M attention). The straight-line surrogate reaches a higher peak (1.00 vs 0.86; 1.00 vs 0.93). \textbf{(b)} Observed $R$ peak (0.81--0.93) against the monotone-warp null (median to p95, 500 draws) and the straight-line surrogate for seven channels: the straight-line surrogate exceeds the observed peak in 7/7; the warp null matches it ($p > 0.05$) in 4/7. The warp null is beaten only in ZIT attention and MLP ($p = 0.014$, $0.008$) and, marginally, ZIB attention ($p = 0.048$). FALSIFIED as a mid-run synchronization event.}
 \label{fig:sync}
\end{figure}

%=====================================================================
\section{An out-of-family check: a convolutional U-Net}\label{sec:unet}

Every model in the archive is a transformer denoiser trained on natural images. To see which results depend on the transformer family, we captured Stable Diffusion~1.5 \citep{rombach2022high}, a latent diffusion model whose denoiser is a convolutional U-Net \citep{ronneberger2015unet} with 16 transformer blocks embedded at four resolutions. Hooks read the self-attention output and the MLP output of each block. Width and token count change with the level: $(N,D)=(4096,320)$, $(1024,640)$, $(256,1280)$ and $(64,1280)$ in the middle block at $512$ pixels. The run used 476 generations on 119 of the archive prompts (about four seeds each), 50 steps, CFG 7.5 and explicit per-branch rows; all numbers below use the conditional branch, and the unconditional branch gives the same verdicts. Three predictions were frozen before the capture existed (SHA-256 prefix \texttt{efa613eb}, 2026-10-08 15:27 UTC), and the operational choices they left open were frozen before scoring (\texttt{7a3b852f}, 21:46 UTC). Every other number in this section is post hoc.

\paragraph{Preregistered.} Because $D$ and $N$ change across blocks, each block is compared with the value $g_{\rm MP}(N,D)$ that a white Gaussian matrix of the same shape gives \citep{marchenko1967distribution}.
\begin{itemize}
 \item \emph{Anisotropy (C1), tested and not supported.} The prediction was $g/g_{\rm MP}<0.10$ in at least $95\%$ of (step, block) cells, as in the transformers, whose median ratios are $0.018$--$0.071$ for attention and $0.0014$--$0.0019$ for the MLP. In SD1.5 the median ratio is $0.077$ for attention, with $31\%$ of cells at or above $0.10$, and $0.26$ for the MLP, with $95\%$ of cells above. The U-Net MLP is about a hundred times closer to isotropy than the transformer MLPs. Strong anisotropy of the uncentred second moment is therefore not a fixed property of denoisers. We do not know the cause; the few extreme activation channels that dominate uncentred transformer spectra \citep{sun2024massive} are one candidate, and this is untested.
 \item \emph{Depth step is not a function of $g$ (C2), confirmed.} With the layer-step definition used throughout, $\beta_{\rm layer}=g_{\ell+1}-g_\ell$, a block effect adds a held-out $\Delta R^2$ of $0.46$ (attention) and $0.61$ (MLP) to a spline in $g$ (5-fold cross-validation grouped by prompt). The same test gives $0.49$--$0.65$ in every transformer channel. In this sense the depth step is non-autonomous in a U-Net as well (Mode-II).
 \item \emph{Depth beyond the level (C3), split.} On pairs of adjacent blocks within one level, block identity adds $\Delta R^2=0.22$ over the level in attention (prediction $\ge0.15$, confirmed) and $0.06$ in the MLP (tested, not supported). In the MLP, the level's width and token count account for almost all of what depth explains.
\end{itemize}

\paragraph{Post hoc.} Each analysis uses the archive method unchanged.
\begin{itemize}
 \item \emph{Backbone.} PC1 is $0.90$ (attention) and $0.62$ (MLP), with effective rank $1.46$ and $2.24$. Both beat all 999 depth-shift nulls (null means $0.22$ and $0.19$; $p=0.001$). The same implementation gives $0.90$ and $0.68$ on SD3.5 Medium and Z-Image Base.
 \item \emph{The spectral triple is two-dimensional.} The first two principal components of the triple $(H,\log Dg,\log\gap)$, with both entropies divided by $\log D$, hold $0.995$ and $0.987$ of the variance, against $0.980$--$0.994$ in the transformers.
 \item \emph{Content geometry.} The ordered depth profile at the first state decodes the ten domains at $0.38$ (attention) and $0.42$ (MLP). This beats every depth scramble (maximum $0.17$ and $0.19$; $p=0.01$) and the layer mean ($0.19$, $0.18$), but only matches the TF-IDF decoder on the prompt text ($0.40$); the transformers exceed it (Table~\ref{tab:decode}). Prompt-relational geometry is shared with every transformer (Fig.~\ref{fig:unet}c): on the 116 prompts common to all six models, Mantel $\rho$ between SD1.5 attention and the five transformers is $0.11$--$0.35$ (every $p=0.0005$), the lower part of the $0.14$--$0.43$ range of unrelated transformer pairs on the same prompts.
 \item \emph{Plateau onset.} SD1.5 uses a variance-preserving $\varepsilon$-prediction schedule. Mapped to the rectified-flow clock at equal signal-to-noise ratio, $t=\sigma/(1+\sigma)$, the onset is $0.58$ (attention) and $0.67$ (MLP), below the archive's $0.78$--$0.87$ at $1024$ pixels. With the SD3 resolution shift for $512$ pixels it is $0.73$ and $0.80$, against $0.84$ for SD3.5 Medium at $512$ pixels on the same correction. The schedule starts at $t=0.94$, so the transient at $\sigma\approx1$ is truncated, and this comparison depends on the clock mapping.
 \item \emph{Prompt share.} Within a sampling state, prompt identity explains $0.56$--$0.67$ of the variance (ANOVA share), inside the archive's $0.36$--$0.95$.
 \item \emph{Clock-axis envelopes.} The layer-mean $g$ of both submodules rises over the whole schedule ($+,+,+$ in every bootstrap, on all three clock mappings; Fig.~\ref{fig:sigma}). The transformers' attention late fall and the Z-Image MLP contraction are both absent. This analysis was defined in writing before it was run (22:45 UTC, \texttt{6251b69f}) but after the data were seen, so it is post hoc.
\end{itemize}
The layer-axis sign structure of \S\ref{sec:topology} was not tested. Pooling $\beta$ across U-Net levels with different $D$ and $N$ is not meaningful, and the preregistration excluded it for that reason.

\begin{figure}[t]
 \centering
 \includegraphics[width=\linewidth]{I-F13_unet}
 \caption{Out-of-family check on the SD1.5 U-Net (conditional branch). \textbf{(a)} Preregistered C1: per-block $g/g_{\rm MP}(N,D)$, median over 476 generations (markers) and range over 50 steps (bars), with the block width above. The transformer archive's median ratios lie in the grey band. The prediction ($<0.10$ in $\ge95\%$ of cells) fails: $31\%$ of attention cells and $95\%$ of MLP cells reach $0.10$. \textbf{(b)} Preregistered C2 and C3: held-out $\Delta R^2$ for $\beta_{\rm layer}$ of a block effect over a spline in $g$ (C2; transformer range in grey) and over the U-Net level on within-level pairs (C3). Dashed: prediction thresholds. \textbf{(c)} Post hoc: Mantel $\rho$ between the prompt-dissimilarity matrices of SD1.5 attention and each transformer on 116 shared prompts at the first, middle and last step (every $p=0.0005$); grey: unrelated transformer pairs on the same prompts.}
 \label{fig:unet}
\end{figure}

What carries over to the U-Net is a low-dimensional backbone, a two-dimensional triple, a depth step that is not a function of $g$, content-bearing depth order and prompt-relational geometry shared with transformers. What does not carry over is the degree of anisotropy, most strongly in the MLP. Domain decoding is weaker, the plateau onset is later on the matched clock, and the clock-axis envelopes rise monotonically. This is one U-Net trained on natural images, so it separates the transformer block from the rest of the setup, but not the data or the objective.

%=====================================================================
\section{An out-of-modality check: an audio flow transformer}\label{sec:ace}

The U-Net changes the block architecture but keeps natural-image data. To change the data modality instead, we captured ACE-Step~1.5 \citep{gong2025acestep}, a music-generation flow transformer with 24 layers and width $D=2048$ whose sampled latent is a 180-second audio clip ($N=2250$ tokens, so $N\ge D$ as in the image captures). Its even layers use sliding-window self-attention over 128 tokens and its odd layers full attention. Hooks read the self-attention and MLP outputs of every layer. The run used the six-step distilled turbo checkpoint at CFG 1 (conditional branch only), 120 generations on 29 instrumental prompts in ten genre domains, mostly four random seeds per prompt. Ten predictions (B1--B10) were frozen before the capture existed, in the same preregistration as the U-Net arm (\texttt{efa613eb}, 2026-10-08 15:27 UTC). The scoring plan, which records how the capture departed from the protocol (turbo checkpoint, six steps, 29 of the 30 prompts) and fixes the operational rules, was frozen before scoring (\texttt{7fa54aab}, 2026-10-09 06:36 UTC). Every number not tied to B1--B10 is post hoc. The instrument gate B1 passed: 24 layers $\times$ two submodules on every step, $N=2250$, $D=2048$.

\paragraph{Preregistered results.}
\begin{itemize}
 \item \emph{Anisotropy (B2), tested and not supported.} The prediction was $g/g_{\rm MP}<0.10$ in at least $95\%$ of (step, layer) cells, with $g_{\rm MP}=0.75$ the isotropic value at this $(N,D)$. The median ratio is $0.15$ for attention and $0.30$ for the MLP, and $72\%$ and $95\%$ of cells reach $0.10$ (Fig.~\ref{fig:ace}a). As in the U-Net, the MLP sits far closer to isotropy than in the image transformers.
 \item \emph{The triple is two-dimensional (B3), tested and not supported.} The first two principal components of $(H/\ln D, \ln(gD)/\ln D, \ln\gap)$ carry $0.988$ of the variance in attention but $0.976$ in the MLP, just below the threshold of $0.98$ (images: $0.980$--$0.994$).
 \item \emph{Depth step is not a function of $g$ (B4), confirmed.} A layer effect adds a held-out $\Delta R^2$ of $0.60$ (attention) and $0.77$ (MLP) to a spline in $g$, against $0.49$--$0.65$ in the image channels.
 \item \emph{Layer parity (B5), confirmed.} Even layers have higher $\log g$ than odd layers in every (generation, step) cell (median difference $0.58$ in attention; the MLP shows the same sign, $0.12$, descriptively). The depth profile is a zig-zag that follows the attention window (Fig.~\ref{fig:ace}a).
 \item \emph{Layer-axis recurrence (B6), tested and not supported.} The analyzer's binned $\beta_{\rm layer}(g)$ for attention has a single crossing ($+,-$ at $g^{*}=0.090$), not the three of the image archive, in 99 of 100 prompt bootstraps (Fig.~\ref{fig:ace}b). The even-layer subsequence alone gives three crossings and the odd-layer subsequence one (descriptive). The crossing-slope prediction B7 depended on B6 and was not run.
 \item \emph{Plateau onset (B8), confirmed.} The onset lies at $\sigma=0.90$ (attention) and $0.80$ (MLP), inside the window $[0.73,0.92]$. With six steps only one interval lies below $\sigma=0.6$, so the plateau speed that defines the onset rests on that interval.
 \item \emph{No susceptibility peak (B9), confirmed.} The seed variance of $\operatorname{logit} g$ has its maximum at the interior step $\sigma=0.60$ in every bootstrap, but it does not beat the last step by two bootstrap standard errors, so no peak passes the preregistered rule (Fig.~\ref{fig:ace}c). The curve is not flat; the rule, not the shape, decides this verdict.
 \item \emph{Prompt share rises first (B10), tested and not supported.} The prompt share of geometry variance is already $0.96$ (attention) and $0.94$ (MLP) at $\sigma=1$, dips to $0.86$ and $0.82$ at $\sigma=0.60$ and recovers at the last step, so its steepest rise is in the last interval, not the first.
\end{itemize}

\paragraph{Post hoc results.} The temporal backbone recurs: PC1 $0.78$ and $0.56$, effective rank $2.10$ and $2.73$, against depth-shift null ranks of $3.9$ and $4.1$ ($p=0.001$ and $0.011$; with six states the rank is at most five, Fig.~\ref{fig:backbone}). The clock-axis envelopes of both submodules peak at $\sigma=1$ and follow $-,+,-$ in every bootstrap (\S\ref{sec:topology}, Fig.~\ref{fig:sigma}). In the analyzer replay (Table~\ref{tab:beta}) attention has a single crossing at all six steps and in every generation bootstrap, and both surrogates reproduce it; the MLP has three crossings at 2 of 6 steps and in $5\%$ of generation bootstraps. Domain decoding is not informative, because 29 prompts give only two or three per domain and even a decoder on the prompt text is at chance (Table~\ref{tab:decode}). Prompt-relational geometry and depth tangents were not compared with the image models, which saw different prompts.

\begin{figure}[t]
 \centering
 \includegraphics[width=\linewidth]{I-F14_ace}
 \caption{Out-of-modality check on the audio flow transformer ACE-Step~1.5 (turbo, 6 steps, cond-only, 120 generations, 29 prompts). \textbf{(a)} Preregistered B2 and B5: per-layer $g/g_{\rm MP}$ ($g_{\rm MP}=0.75$), median over generations (markers) and range over steps (bars); even layers (sliding-window attention) filled, odd layers (full attention) open. Grey band: image-transformer median ratios. B2 fails ($72\%$ of attention and $95\%$ of MLP cells reach $0.10$); the even--odd zig-zag holds in every cell (B5). \textbf{(b)} Preregistered B6: the analyzer's binned $\beta_{\rm layer}(g)$, all steps pooled. Attention and MLP each cross zero once, not three times; the even-layer subsequence alone crosses three times (descriptive). \textbf{(c)} Preregistered B9: seed variance of $\operatorname{logit} g$ along $\sigma$. The interior maximum at $\sigma=0.60$ does not clear the last step by two bootstrap SE, so no peak passes the rule.}
 \label{fig:ace}
\end{figure}

What carries over to audio is the low-dimensional backbone, a depth step that is not a function of $g$, and a plateau onset inside the image range. What does not carry over is the strong anisotropy, the depth sign sequence $+,-,+,-$, the early rise of prompt share and, narrowly, the two-dimensional triple in the MLP. By the consequence fixed in the preregistration, the layer-axis crossing pattern is specific to the image and video data of this corpus, and the recurrence of coarse organisation along depth is scoped to those modalities. This is a single model, and it is distilled, sampled in six steps and captured at CFG 1; Z-Image Turbo attention, also distilled and conditional-only, gives a single crossing too (Table~\ref{tab:beta}). An undistilled audio model with more steps would separate modality from distillation.

%=====================================================================
\section{Synthesis}\label{sec:synthesis}
\begin{enumerate}
 \item Coarse organisation recurs across image and video architectures: a low-dimensional backbone, an early transient followed by a $\sigma$-plateau, and the depth sign sequence $+,-,+,-$ at some noise level in every model.
 \item Prompt-relational geometry is broadly shared, beyond theme and text, and strongest in the lineage.
 \item Local depth dynamics are lineage-specific: tangents separate the distillation pair from every unrelated pair by $0.70$.
 \item The same data therefore support a hierarchy of invariance, not one universal microscopic law.
 \item In a convolutional U-Net the backbone, the non-autonomous depth step and shared prompt geometry recur, but the strong anisotropy of the transformer channels and their clock-axis envelope shapes do not (\S\ref{sec:unet}).
 \item In an audio flow transformer the backbone, the non-autonomous depth step and the plateau onset recur, but the strong anisotropy, the depth sign sequence and the early rise of prompt share do not (\S\ref{sec:ace}). By the preregistered consequence, the depth sign sequence is specific to image and video data in this corpus.
\end{enumerate}

%=====================================================================
\section{Limitations and future tests}\label{sec:limits}
\begin{itemize}
 \item \textbf{Commitment.} A separatrix or commitment claim needs cloned continuations around candidate noise levels with scored final images (gate G3); recent work on semantic commitment in diffusion \citep{handke2026entropic,raya2023spontaneous} motivates it.
 \item \textbf{Autoregressive dynamics.} Per-token phases need context-invariant statistics (G4).
 \item \textbf{Output quality.} The archive has no stable output identifier or quality label, so no geometry$\to$quality statement is possible (D14).
 \item \textbf{Response.} Finite perturbations would turn the descriptive local maps into response or Jacobian measurements (G1).
 \item \textbf{Submodules.} Among the transformers, MLP hooks exist only for Z-Image; the clock-axis result needs them on Mochi and SD3.5. The U-Net and the audio model have MLP hooks, and neither MLP shows the Z-Image clock-axis pattern (post hoc, \S\ref{sec:topology}).
 \item \textbf{Massive activations.} Uncentred spectra can be dominated by a few extreme channels \citep{sun2024massive,sun2026anatomy}, which in DiTs are tied to the timestep modulation \citep{gan2025massive,gan2025unleashing}. Recomputing with those channels removed or centred is a needed control; because they also carry prompt information \citep{turri2026few}, removal may remove content signal as well.
 \item \textbf{Origin of the recurring structure.} We make no claim about where the coarse organisation comes from. Every transformer channel here shares natural-image statistics, a denoising objective whose noise schedule already orders generation coarse-to-fine in frequency \citep{falck2025fourier}, and residual transformer blocks. The U-Net check (\S\ref{sec:unet}) changes the block architecture only, and leaves the data and the objective unchanged. The audio check (\S\ref{sec:ace}) changes the data modality and keeps a transformer and a flow-matching objective; it is a single distilled six-step model, so modality is confounded with distillation and step count; models may also converge on shared statistics of their inputs \citep{huh2024platonic} without abstracting the underlying laws \citep{kang2025videophysics}. Separating inherited data statistics from a more general regularity needs matched models trained on data with altered statistics, a different objective or a different architecture, and context-invariant language-model captures (G4).
 \item \textbf{Capture protocol.} Two attempts to lock a frozen prompt file and a unique-seed design failed on real captures; future captures should use the shipped prompt file and preflight.
\end{itemize}

\section*{Code and data}
All scripts, figure code (Python) and derived tables are at \url{https://github.com/Z-L-D/mode-I}.

\bibliographystyle{plainnat}
\bibliography{references}

\appendix
\section{Capture and provenance}\label{app:provenance}
The passive archive (\texttt{TELEMETRY.zip}, SHA-256 prefix \texttt{d7f5a82e}, 506{,}345{,}966 bytes) holds the six diffusion campaigns of Table~\ref{tab:corpus} and the language-model controls; one further Z-Image Base domain file (48 generations) was supplied separately, bringing Base to 476 generations on 117 prompts. Generation identifiers are namespaced by campaign and seed. The Round 3 and Round 4 captures were made with an updated capture node that writes per-branch rows, a prompt-influence column and the seed; a repeated CFG 1 capture reproduces the earlier one in every cell to six decimals, and the synthetic branch mean equals the per-branch mean to $2\times10^{-16}$. The Turbo archive also contains four auxiliary 100-step generations, used only for the step-count control.

\section{Preregistration scorecards}\label{app:prereg}
\begin{table}[h]
\centering\small
\begin{tabular}{llllc}
\toprule
Round & SHA-256 & Frozen (UTC) & Rows & Confirmed / partial / refuted\\
\midrule
1 (review) & -- & 2026-10-05 & -- & not preregistered\\
2 & \texttt{2343fd02} & 2026-10-05 11:05 & 28 & 18 / 3 / 7\\
3 & \texttt{d29fdfa6} & 2026-10-05 15:21 & 32 & 1 / -- / 2 (29 not scorable)\\
4 & \texttt{82ef210a} & 2026-10-06 23:01 & 16 & 11 / -- / 5\\
4 (captures) & \texttt{82ef210a} & 2026-10-06 23:01 & 12 & 5 / -- / 1 (5 not scorable)\\
5 & \texttt{78b250a2} & 2026-10-08 05:24 & 22 & 13 / -- / 9\\
6 & \texttt{c0279f03} & 2026-10-08 07:06 & 22 & 17 / -- / 5\\
7 (release audit) & \texttt{9012d20f} & 2026-10-08 09:02 & 8 & consistency checks only\\
R3 Part C (U-Net) & \texttt{efa613eb} & 2026-10-08 15:27 & 3 & 1 / 1 / 1 (C3 split by channel)\\
R3 Part B (ACE-Step) & \texttt{efa613eb} & 2026-10-08 15:27 & 10 & 4 / -- / 4 (B1 gate; B7 not run)\\
8 (correlations) & \texttt{acc578ac} & 2026-10-09 08:30 & 16 & 7 / -- / 9\\
9 (salvage) & \texttt{09c1f170} & 2026-10-09 09:07 & 12 & 6 / -- / 6\\
10 (final sweep) & \texttt{2338e502} & 2026-10-09 22:27 & 2 & 0 / -- / 2\\
\bottomrule
\end{tabular}
\end{table}
Round 3 predictions were mostly unscorable because the captures did not match the frozen prompt file. In the U-Net arm, C1 (per-block anisotropy) is refuted in both channels, C2 (depth step not a function of $g$) is confirmed in both, and C3 (depth beyond the level) is confirmed in attention and refuted in the MLP (\S\ref{sec:unet}); its scoring plan was frozen separately (\texttt{7a3b852f}, 21:46 UTC) because the preregistration left the cell definition and the C3 comparison open. In the audio arm (\S\ref{sec:ace}), B4, B5, B8 and B9 are confirmed and B2, B3, B6 and B10 refuted; B7 was conditional on B6 and was not run. Its scoring plan, which records the deviations of the capture from the protocol, was frozen before scoring (\texttt{7fa54aab}, 2026-10-09 06:36 UTC). Rounds 3--5 refuted several of the project's own earlier readings: the synchronisation pulse, the $\sigma_{50}$ onset contrast, the historical scalar ordering and an apparent pooled-capture seed collapse.

\section{Analyzer instrument notes}\label{app:analyzer}
The layer-axis results use the CoS-PAFS Live Analyzer v3 (SHA-256 prefix \texttt{4089062551a3c900}) without modification. Four behaviours matter when it is applied to diffusion telemetry:
\begin{itemize}
 \item Its mean-profile, interval and ratio tools keep only rows with token type \texttt{generated}, so they return nothing on diffusion rows (type \texttt{step}); they were built for the language-model corpus. The layer-flow measurement on diffusion data is the dashboard's $\beta$-layer view.
 \item When a file contains both attention and MLP rows, most views average them per (layer, step) cell. Pooled Z-Image channels then show a single crossing.
 \item A header-less CSV is silently dropped by the loader. One Mochi CFG 4.5 domain file (73{,}152 rows) is affected, so Mochi CFG 4.5 results cover 9 of 10 domains.
 \item The Turbo archive includes four auxiliary 100-step generations, which enter step-axis views up to step 99.
\end{itemize}

\section{Estimator sensitivity}\label{app:estimators}
\begin{table}[h]
\centering\small
\caption{$\sigma_{50}$ and relative onset against sample size (SD3.5 Medium archive, branch-mean).}
\begin{tabular}{lcccc}
\toprule
Prompts $\times$ seeds & $n$ & $\sigma_{50}$ & Relative onset & Final $R^2$\\
\midrule
$30\times2$ & 60 & undefined (98\%) & 0.900 & 0.48\\
$60\times2$ & 120 & 0.720 & 0.867 & \\
$30\times4$ & 120 & 0.867 & 0.935 & \\
$120\times2$ & 240 & 0.843 & 0.915 & \\
$60\times4$ & 240 & 0.917 & 0.952 & \\
$120\times4$ & 480 & 0.936 & 0.956 & 0.89\\
\bottomrule
\end{tabular}
\end{table}
The fold partition alone moves $\sigma_{50}$ between $0.48$ and $0.63$ on one dataset. Simulated references for VR(4) on paths of the archive's length: random walk $0.816$--$0.828$, AR(0.8) $0.694$--$0.725$, i.i.d.\ $0.25$.

\section{Legacy instrument results}\label{app:legacy}
An early SD3.5 Medium suite (one prompt, 30 seeds per resolution, CFG 1, 100 steps, $512/768/1024$ pixels) and a CFG 1.0 versus 5.5 study at $256$ pixels were recorded with a capture path that substitutes an exponential-moving-average attention whenever it is registered. On those files depth profiles show a terminal rebound with four crossings, shape correlations of $0.989$--$0.996$ across resolutions, and a stable complex eigenpair of the local depth map. Because the instrument differs and the suite has a single prompt, these results are not used as evidence in this paper. The Round 3 captures at the same resolutions, with 60 prompts on the patched instrument, show $+,-,+,-$ rather than four crossings (\S\ref{sec:interventions}).

\section{Supplementary figures}\label{app:supp}
\renewcommand{\thefigure}{S\arabic{figure}}\setcounter{figure}{0}
\begin{figure}[h]
 \centering
 \includegraphics[width=\linewidth]{I-S1_theme_residual}
 \caption{Same-theme residual alignment after each model/time backbone is removed, for all model pairs at the first, middle and last step. Bars show the theme-label permutation null p95. Base$\leftrightarrow$Turbo: 0.63--0.75 at every step, the strongest; At the first step every other pair is $\le 0.24$; By the middle step unrelated pairs reach up to 0.48, and up to 0.52 by the last step (ZIT$\leftrightarrow$SD3.5\,M). Theme alignment is strongest in the lineage at every step and lineage-specific only at the first step.}
 \label{fig:S1}
\end{figure}

\begin{figure}[h]
 \centering
 \includegraphics[width=\linewidth]{I-S2_stepcount_controls}
 \caption{Per-$\sigma$ geometry changes by at most 2--3\% between 8/40 and 100 steps. \textbf{(a)} SD3.5\,M, one prompt, explicit cond branch: layer-mean attention $g(\sigma)$ at 40 (markers) and 100 steps (lines), CFG\,1 and CFG\,4.5. The same-prompt 40$\to$100 shift is 0.033 at CFG\,1 (seed-split null p95 0.027, $p = 0.023$) and 0.007 at CFG\,4.5 (p95 0.022, $p = 0.70$). \textbf{(b)} Mochi 30 vs 100 steps: path speed by $\sigma$ band (ratio 0.81--1.06). The two arms also differ in CFG (4.5 vs 1). \textbf{(c)} ZIT 8 vs 100 steps at matched $\sigma$: median relative difference of $g$ (1.9\% attn / 1.4\% MLP), gap (1.4\% / 0.4\%) and $H$ (0.2\% / 0.5\%). The 100-step set contains only 4 generations.}
 \label{fig:S2}
\end{figure}

\begin{figure}[h]
 \centering
 \includegraphics[width=\linewidth]{I-S3_qwen_context_clock}
 \caption{Why per-token phase claims for autoregressive models are not made. The v1 LLM capture re-encodes the full context at every token without a KV cache, so each logged value is a running statistic of a nested $n\times D$ matrix. \textbf{(a, b)} Depth-mean $g$ and $H$ of generated tokens against sequence length $n$ for Qwen2.5 1.5B/3B/7B (median per $n$ bin). \textbf{(c)} $g$ normalized by the finite-sample baseline $n/(n+D)$ remains monotone in $n$. \textbf{(d)} Median within-generation Spearman $\rho(n,\cdot)$ for all seven LLM captures: $\rho = 1.00$ for $g$ and $\le -0.999$ for PR$/n$ and $\tilde g$. The bottom row is a model-free reference (Round 9, preregistered): nested running statistics of i.i.d.\ Gaussian rows (180 simulations; isotropic, power-law and spiked spectra; $D=1536$ and $4096$) give the same $\rho$ without any network, so the monotonicity is a property of the capture. The Qwen multiscale $\beta$ "anomaly" therefore is a context-length clock, and it is demoted.}
 \label{fig:S3}
\end{figure}

\begin{figure}[h]
 \centering
 \includegraphics[width=\linewidth]{I-S4_negative_atlas}
 \caption{Retired claims. Every panel is marked RETIRED. \textbf{(a)} "The semantic decoder maximum is the universal C1 time": the smoothed semantic Fisher peak (0.14--0.69) does not track the A-speed 20\% onset (0.30--0.79) in normalized time (differences $-0.21$ to $+0.40$). \textbf{(b)} "Trajectory curvature gives a universal C1 detector": the time of the maximum turning angle falls in the early transient ($\tau \le 0.14$) in six of seven channels and at $\tau = 0.96$ in SD3.5\,L. \textbf{(c)} "Universal seed collapse": seed share of within-step variance along $\sigma$ (ANOVA). It ranges from 0.05--0.08 (ZIT attention) to a peak of 0.64 (ZIB MLP), with no common collapse. The retired passive-susceptibility peak is not shown, because no per-channel values are stored.}
 \label{fig:S4}
\end{figure}

\end{document}
